\documentclass[12pt]{article}
\usepackage[a4paper,margin=1in]{geometry}
\usepackage[utf8]{inputenc}
\usepackage[T1]{fontenc}
\usepackage{amsmath,amssymb}
\usepackage{textcomp}
\usepackage{parskip}
\begin{document}
\begin{center}
{\Large \textbf{1998 STEP II}}
\end{center}

\subsection*{Page 1}
\begin{flushleft}
STEP II, 1998\par
\bigskip\textbf{Section A: Pure Mathematics}\par
\bigskip\textbf{Question 1.} Show that, if n is an integer such that\par
3 3 3\par
\[(n 3) + n = (n + 3) , ( )\]
- ∗\par
then n is even and n2 is a factor of 54. Deduce that there is no integer n which satisfies the\par
equation ( ).\par
∗\par
Show that, if n is an integer such that\par
3 3 3\par
\[(n 6) + n = (n + 6) , ( )\]
- ∗∗\par
then n is even. Deduce that there is no integer n which satisfies the equation ( ).\par
∗∗\par
\bigskip\textbf{Question 2.} Use the first four terms of the binomial expansion of (1 1/50)1/2, writing 1/50 = 2/100 to\par
-\par
\[simplify the calculation, to derive the approximation \sqrt\{2\} 1.414214.\]
\approx\par
Calculate similarly an approximation to the cube root of 2 to six decimal places by considering\par
(1 + N/125)a, where a and N are suitable numbers.\par
[You need not justify the accuracy of your approximations.]\par
\bigskip\textbf{Question 3.} Show that the sum S of the first N terms of the series\par
N\par
1 3 5 2n 1\par
+ + + + - +\par
1.2.3 2.3.4 3.4.5 \cdots n(n + 1)(n + 2) \cdots\par
is\par
1 3 1 5\par
+ .\par
2 µ 2 N + 1 - N + 2 ¶\par
What is the limit of S as N ?\par
N\par
\to \infty\par
The numbers a are such that\par
n\par
a (n 1)(2n 1)\par
n\par
\[= - - .\]
a n-1 (n + 2)(2n 3)\par
-\par
\infty 2\par
Find an expression for a\par
n\par
\[/a^1 and hence, or otherwise, evaluate a\]
n\par
\[when a^1 = .\]
\[nP =1\]
\end{flushleft}
\subsection*{Page 2}
\begin{flushleft}
\bigskip\textbf{Question 4.} The integral I is defined by\par
n\par
\[\pi\]
\[I = (\pi/2 x) sin(nx + x/2) cosec (x/2) dx,\]
n\par
\[Z^0 -\]
where n is a positive integer. Evaluate I\par
n\par
I n-1, and hence evaluate I\par
n\par
leaving your answer\par
-\par
in the form of a sum.\par
\[5 Define the modulus of a complex number z and give the geometric interpretation of z^1 z^2\]
| - |\par
for two complex numbers z1 and z2. On the basis of this interpretation establish the inequality\par
\[z^1 + z^2 6 z^1 + z^2 .\]
| | | | | |\par
\[Use this result to prove, by induction, the corresponding inequality for z^1 + + z\]
n\par
.\par
| \cdots |\par
\[The complex numbers a^1, a^2, . . . , a\]
n\par
satisfy a\par
i\par
\[6 3 (i = 1, 2, . . . , n). Prove that the equation\]
| |\par
\[a1z + a2z 2 + a\]
n\par
\[zn = 1\]
\cdots\par
\[has no solution z with z 6 1/4.\]
| |\par
\bigskip\textbf{Question 6.} Two curves are given parametrically by\par
\[(1) x^1 = (θ + sin θ), y^1 = (1 + cos θ),\]
and\par
\[(2) x^2 = (θ sin θ), y^2 = (1 + cos θ).\]
- -\par
Find the gradients of the tangents to the curves at the points where θ = \pi/2 and θ = 3\pi/2.\par
\[Sketch, using the same axes, the curves for 0 6 θ 6 2\pi.\]
Find the equation of the normal to the curve (1) at the point with parameter θ. Show that\par
this normal is a tangent to the curve (2).\par
\end{flushleft}
\subsection*{Page 3}
\begin{flushleft}
7 Let\par
\[f(x) = tan x x,\]
-\par
\[g(x) = 2 2 cos x x sin x,\]
- -\par
\[h(x) = 2x + x cos 2x sin 2x,\]
- 2\par
\[x(cos x)1/3\]
\[F(x) = .\]
sin x\par
\textbf{(i)} By considering f(0) and f (x), show that f(x) > 0 for 0 < x < \pi/2.\par
\textbf{(ii)} Show similarly that g(x) > 0 for 0 < x < \pi/2.\par
\textbf{(iii)} Show that h(x) > 0 for 0 < x < \pi/4, and hence that\par
2 2\par
\[x(sin x + 3 cos x) 3 sin x cos x > 0\]
-\par
\[for 0 < x < \pi/4.\]
F (x)\par
\textbf{(iv)} By considering , show that F (x) < 0 for 0 < x < \pi/4.\par
F(x)\par
\bigskip\textbf{Question 8.} Points A, B, C in three dimensions have coordinate vectors a, b, c, respectively. Show that\par
the lines joining the vertices of the triangle ABC to the mid-points of the opposite sides meet\par
at a point R.\par
P is a point which is not in the plane ABC. Lines are drawn through the mid-points of BC,\par
CA and AB parallel to P A, P B and P C respectively. Write down the vector equations of the\par
lines and show by inspection that these lines meet at a common point Q.\par
Prove further that the line P Q meets the plane ABC at R.\par
\end{flushleft}
\subsection*{Page 4}
\begin{flushleft}
\bigskip\textbf{Section B: Mechanics}\par
\bigskip\textbf{Question 9.} A light smoothly jointed planar framework in the form of a regular hexagon ABCDEF is\par
suspended smoothly from A and a weight 1kg is suspended from C. The framework is kept\par
rigid by three light rods BD, BE and BF . What is the direction and magnitude of the\par
supporting force which must be exerted on the framework at A?\par
Indicate on a labelled diagram which rods are in thrust (compression) and which are in tension.\par
Find the magnitude of the force in BE.\par
\bigskip\textbf{Question 10.} A wedge of mass M rests on a smooth horizontal surface. The face of the wedge is a smooth\par
plane inclined at an angle α to the horizontal. A particle of mass m slides down the face of\par
the wedge, starting from rest. At a later time t, the speed V of the wedge, the speed v of the\par
particle and the angle β of the velocity of the particle below the horizontal are as shown in\par
the diagram.\par
Let y be the vertical distance descended by the particle. Derive the following results, stating\par
in (ii) and (iii) the mechanical principles you use:\par
\textbf{(i)} V sin α = v sin(β α);\par
-\par
\textbf{(ii)} tan β = (1 + m/M ) tan α;\par
\textbf{(iii)} 2gy = v2(M + m cos2 β)/M .\par
Write down a differential equation for y and hence show that\par
\[gM t^2 sin 2 β\]
\[y = .\]
\[2 (M + m cos2 β)\]
\end{flushleft}
\subsection*{Page 5}
\begin{flushleft}
\bigskip\textbf{Question 11.} A fielder, who is perfectly placed to catch a ball struck by the batsman in a game of cricket,\par
watches the ball in flight. Assuming that the ball is struck at the fielder's eye level and is\par
d\par
caught just in front of her eye, show that (tan θ) is constant, where θ is the angle between\par
dt\par
the horizontal and the fielder's line of sight.\par
In order to catch the next ball, which is also struck towards her but at a different velocity, the\par
fielder runs at constant speed v towards the batsman. Assuming that the ground is horizontal,\par
d\par
show that the fielder should choose v so that (tan θ) remains constant.\par
dt\par
\end{flushleft}
\subsection*{Page 6}
\begin{flushleft}
\bigskip\textbf{Section C: Probability and Statistics}\par
\bigskip\textbf{Question 12.} The diagnostic test AL has a probability 0.9 of giving a positive result when applied to a\par
person suffering from the rare disease mathematitis. It also has a probability 1/11 of giving a\par
false positive result when applied to a non-sufferer. It is known that only 1\% of the population\par
suffer from the disease. Given that the test AL is positive when applied to Frankie, who is\par
chosen at random from the population, what is the probability that Frankie is a sufferer?\par
In an attempt to identify sufferers more accurately, a second diagnostic test STEP is given\par
to those for whom the test AL gave a positive result. The probablility of STEP giving a\par
positive result on a sufferer is 0.9, and the probability that it gives a false positive result on a\par
non-sufferer is p. Half of those for whom AL was positive and on whom STEP then also gives\par
a positive result are sufferers. Find p.\par
\bigskip\textbf{Question 13.} A random variable X has the probability density function\par
\[λe -λx x > 0,\]
\[f (x) =\]
\[½ 0 x < 0.\]
Show that\par
\[P(X > s + t X > t) = P(X > s).\]
|\par
The time it takes an assistant to serve a customer in a certain shop is a random variable with\par
the above distribution and the times for different customers are independent. If, when I enter\par
the shop, the only two assistants are serving one customer each, what is the probability that\par
these customers are both still being served at time t after I arrive?\par
One of the assistants finishes serving his customer and immediately starts serving me. What\par
is the probability that I am still being served when the other customer has finished being\par
served?\par
\bigskip\textbf{Question 14.} The staff of Catastrophe College are paid a salary of A pounds per year. With a Teaching\par
Assessment Exercise impending it is decided to try to lower the student failure rate by offering\par
each lecturer an alternative salary of B/(1 + X) pounds, where X is the number of his or\par
her students who fail the end of year examination. Dr Doom has N students, each with\par
independent probability p of failure. Show that she should accept the new salary scheme if\par
\[A(N + 1)p < B(1 (1\]
p)N+1\par
).\par
- -\par
Under what circumstances could X, for Dr Doom, be modelled by a Poisson random variable?\par
What would Dr Doom's expected salary be under this model?\par
\end{flushleft}
\end{document}
